\documentclass[twocolumn]{aastex631}
\begin{document}
\title{An age dependence of the Galactic alpha-knee across galactocentric radius}
\author{NebulaMind Lab (autonomous pipeline)}
\affiliation{NebulaMind Lab, \url{https://lab.nebulamind.net}}
\begin{abstract}
Age-resolved alpha-knee from an APOGEE (DR18) 3-table join of 330,446 giants (apogeeDistMass C/N ages + distances, apogeeStar Galactic coordinates, aspcapStar [Fe/H]-[SI/Fe]). The [Fe/H]\_knee is located per (R\_g x age) cell with a broken-line ridge fit (bootstrap error) in 31 populated cells; median [Fe/H]\_knee = -0.45. Gradients: d[Fe/H]\_knee/d(age)|\_R\_g = +0.0506+/-0.0093 dex/Gyr; d[Fe/H]\_knee/dR\_g|\_age = +0.0306+/-0.0086 dex/kpc. Non-circularity: the age-gradient sign HOLDS under the abundance-scale swap ([Fe/H]->[M/H], [SI/Fe]->[alpha/M]). Generated autonomously from public data (apogee) via the NebulaMind Lab runner: a bounded, reproducible, descriptive study.
\end{abstract}
\section{Introduction}
The age-resolved alpha-knee has been a topic of interest in understanding the chemical evolution of the Milky Way. Previous studies have explored various methods for determining stellar ages and their relationship with chemical abundances [Claytor2020, Valle2024]. However, these works primarily focused on specific populations or lacked comprehensive spatial resolution. Grisoni et al. (2024) highlighted the importance of considering guiding radius and Galactic height in studying alpha-rich stars, emphasizing the need for a more detailed analysis of the alpha-knee's dependence on age and location [Grisoni2024]. Building upon these findings, our research aims to investigate the relationship between the alpha-knee, stellar age, and galactic radius using APOGEE data.
\section{Data and method}
To achieve this goal, we utilized raw SDSS DR18 APOGEE catalog data from SkyServer, extracting information from three tables: apogeeDistMass, apogeeStar, and aspcapStar. These tables were joined in-process on APOGEE\_ID after applying flag/quality cuts to ensure data reliability. We computed R\_g values using glon/glat/distance with R0=8.122 kpc and derived ages from spectroscopic C/N ratios. The alpha-knee was determined through a broken-line ridge fit, accounting for bootstrap errors in 31 populated cells.
\section{Result}
Our analysis yielded an age-resolved alpha-knee measurement: the median [Fe/H]\_knee value is -0.45, with gradients of d[Fe/H]\_knee/d(age)|\_R\_g = +0.0506+/-0.0093 dex/Gyr and d[Fe/H]\_knee/dR\_g|\_age = +0.0306+/-0.0086 dex/kpc. Notably, the age-gradient sign remains consistent even when swapping the abundance calibration scale ([Fe/H]->[M/H], [SI/Fe]->[alpha/M]).
\begin{figure}\centering\includegraphics[width=\columnwidth]{result.png}\caption{An age dependence of the Galactic alpha-knee across galactocentric radius}\end{figure}
\section{Caveats}
Despite these findings, it is essential to acknowledge the limitations of our approach. The automated selection process may introduce biases, as it relies on a single method for determining stellar ages and does not account for potential systematic errors in the C/N-age relationship. Additionally, the uncalibrated nature of our measurement means that absolute age values should be treated with caution. Furthermore, our analysis assumes a fixed R0 value, which might not accurately represent the complex geometry of the Milky Way. These caveats highlight the need for further refinement and validation of our results using complementary methods and data sets.
\section*{References}
\footnotesize
[Claytor2020] Claytor et al. (2020). Chemical Evolution in the Milky Way: Rotation-based Ages for APOGEE-Kepler Cool Dwarf Stars. bibcode 2020ApJ...888...43C \par [Grisoni2024] Grisoni et al. (2024). K2 results for "young" α-rich stars in the Galaxy. bibcode 2024A\&A...683A.111G \par [Valle2024] Valle et al. (2024). Stellar model tests and age determination for RGB stars from the APO-K2 catalogue. bibcode 2024A\&A...690A.323V \par [Warfield2021] Warfield et al. (2021). An Intermediate-age Alpha-rich Galactic Population in K2. bibcode 2021AJ....161..100W

\end{document}
